% !TeX root = ../papers/nc-p3.tex
\section{被約な点を持つ有限窓の行列交換子と三次表示}
\label{sec:marked-point-commutators}
有限Frobeniusの四辺が有限で、長さ1の局所成分を一つ指定する。
前節の四周期列を用いてその評価を各位置で $t_i$ に1、補空間 $H_i$ で0とする。
\begin{theorem}\label{thm:marked-point-commutators}
印付きテンソルは二つの三次元空間 $U,V$ と
\[
 G\in U\otimes\mathfrak{sl}(V),\quad F\in U^*\otimes\mathfrak{sl}(V),\quad
 Q\in V\otimes\mathfrak{sl}(U),\quad P\in V^*\otimes\mathfrak{sl}(U)
\]
によって次の形にできる。位置 $0,1,2,3$ の $H_i$ を $U,V,U^*,V^*$ とし、
各補文字の添字を $*$、残りを $a,b,c,d=1,2,3$ とする。
隣接する二つの $*$ を持つ係数は零で、他の係数は
\[
\begin{array}{ll}
 w_{*b*d}=\delta_{bd},&w_{a*c*}=\delta_{ac},\\
 w_{*bcd}=(F_c)_{bd},& w_{ab*d}=(G_a)_{bd},\\
 w_{a*cd}=(P_d)_{ac},& w_{abc*}=(Q_b)_{ac},\\
 w_{abcd}=(G_a F_c)_{bd}+(Q_b P_d)_{ac}.&
\end{array}
\]
二つの隣接切断の階数6はこの表示ではちょうど
\[
 [G_a,F_c]_{bd}+[Q_b,P_d]_{ac}=0\qquad(1\le a,b,c,d\le3)
\]
に等しい。各三行列組は三次元空間に共通の不変直線も不変平面も持たない。
印を保つ同一視は $A\in\mathrm{GL}(U), B\in\mathrm{GL}(V), s\in k^*$ による
四因子 $(A,B,s A^{-t},s B^{-t})$ の作用と全体の $s^{-1}$ 倍で尽きる。
$*$ 座標はこの作用で固定する。
\end{theorem}
\begin{proof}
隣接する二点での縮約は零で、被約点における反対位置の二つの行列は前節より階数3である。
独立な基底変更とテンソルの規格化で両行列を単位行列にする。
$t_i$ の $H_i$ による持上げ変更は、反対位置の $*$ 一個を持つ三係数行列に
それぞれ独立にスカラー行列を加える。他の三組には影響しないため、
標数零で跡零の持上げが一意に取れる。
$01\mid23$ 切断のうち各側で一つの $*$ を持つ六行六列は置換単位行列である。
従って階数6なら残る九行九列はそのSchur補元を零にする唯一の表示で、
これが上の $w_{abcd}$ である。他の切断で同じことを行うと二積の順序がそれぞれ逆になる。
よって両条件の差が表示の交換子である。各因子の射影スカラーで $*$ の倍率を1に戻せる。
二対合を同じスカラーまで保つことと跡零の一意性で残余群を得る。
逆にこの表で作った係数について、三次・四次重なりと有限点条件は元の有限行列検査で判定できる。
正則性はこの表の二切断条件とは別である。

例えば $F_c$ が共通の不変直線を持つと、$I,F_1,F_2,F_3$ をその直線のベクトルで評価した像は
一次元であり、$w(t_0^*)$ の端からの縮約が可分解な二次関係になる。
双対側の直線にも同じ議論が適用できるので共通不変平面もない。ほかの三組も同様である。
\end{proof}

\begin{corollary}\label{cor:point-commutator-rank}
未認証の場合（全ての中支持が4）、$G,F,Q,P$ は各々三独立行列を与える。
反対階数は、$g=\dim\langle G_a,F_c\rangle$, $q=\dim\langle Q_b,P_d\rangle$ として
\[
 \rho=2+g+q+\operatorname{rank}\bar w .
\]
ここで $\bar w$ は九行九列の上の $w_{abcd}$ を反対切断 $ac\mid bd$ で読み、
それぞれ $\langle I,Q_b,P_d\rangle$, $\langle I,G_a,F_c\rangle$ を法として得るテンソルである。
特に反対階数8では $G,F$ は同じ三平面の二基底、$Q,P$ も同じ三平面の二基底であり $\bar w=0$ である。
\end{corollary}
\begin{proof}
例えば $w(t_0^*)$ の中央評価空間は $I,F_1,F_2,F_3$ の包である。
反対切断で少なくとも一つ $*$ を含む行は $I,G_a,F_c$ を残りの九列にのみ持つ。
列側でも対称的に $I,Q_b,P_d$ を持ち、印のある行と列の間は零である。
行と列を消去して表示式が従う。
\end{proof}

\begin{proposition}\label{prop:point-kernel-cubic-presentation}
指定した四周期点への評価の核は、三文字 $H_i$ だけを使う正次数の代数である。
この代数に各頂点の単位を加えたものを $B$ とすると、関係は三次だけであり、各位置に八個ある。
たとえば $i=0$ では $H_0\otimes H_2=\operatorname{End}(U)$ と読み、
三次関係の空間は
\[
 Z\longmapsto \big([Q_b,Z]_{ac}\big)_{abc}
 \qquad (Z\in\mathfrak{sl}(U))
\]
の像である。他の位置は巡回的に得る。$B$ の三文字点対応は、中央点で評価した行列と、
両端点の与える階数1の行列が交換する条件である。
各正の長さで $A_{i,j}=B_{i,j}\oplus k\,t_i\cdots t_{j-1}$ であり、
これは $B$ の三次表示から全次数を復元する。
\end{proposition}
\begin{proof}
前節の六関係をここでは $t h=L(h),\ h t=N(h)$ と書き、位置は積順で補う。
純 $H$ 語上で $L$ は最初の文字に、$N$ は最後の文字にこの置換を行う作用素とする。
三文字 $a t b$ の二つの書換えの差は $N(a)b-aL(b)$ である。
表から $N_0(a)=-\sum_b Q_b a\otimes e_b$,
$L_1(c)=-\sum_b e_b\otimes Q_b^t c$ となるので、その差の像は表示の交換子空間である。
三行列組の既約性とSchurの補題により共通中心化環はスカラーであり、像の次元は8である。
他の位置も同様である。点対応の条件はこの写像を三点で評価し、
全 $Z$ に対するtrace対合の非退化性を用いれば従う。

二次書換えで三次の経路空間は全 $t$ の一語と純 $H$ の27語で張られ、
有限Frobenius検査からその次元は20である。従って純 $H$ の三次関係は既に得た八個で尽きる。
特に $t h t$ の差 $LN(h)-NL(h)$ もそれらの包にある。
この八関係だけによる自由 $H$ 代数の商で $L,N$ が良定義であることを確かめる。
$N(a)b-aL(b)$ の左端に $L$ を施すと、
$LN(a)b-L(a)L(b)$ となる。新しい語の中央に同じ八関係を使えば
$L(a)L(b)=NL(a)b$ となるため零である。右端での $N$ も逆順で同様である。
関係が端から離れていれば置換はそのまま前後の語に作用する。よって生成イデアルは保たれる。
両置換は長さ2以上で別の文字に作用するので交換し、長さ1でも先の $t h t$ の式で交換する。
さらに任意の非空な純語で $N(a)b=aL(b)$（隣り合う末尾・先頭に八関係を適用）である。

従ってこの三次商の正次数部分と $k[t]$ の反復の直和に、
左 $t$ を $L$、右 $t$ を $N$ とする結合積を入れられる。
結合律は $L(ab)=L(a)b$, $N(ab)=aN(b)$、および直前の二等式である。
この直和が元の六つの二次式を満たし、逆にそれらから全ての指定積が従うので
互いに逆な準同型を得る。
\end{proof}
